% Draft: domain-specific exclusion radius for genomic negative controls
% Target: Chapter 4 (negative construction); one sentence in Chapter 7.
% Status: draft v1, 2026-11-02. References checked against Crossref DOIs.
% Project sources: log R-036, specs/graphpro_dataset_reconstruction.md,
%   specs/sequence_window_policy.md, F-117.

\section{Exclusion Radius for Genomic Negative Controls}\label{draft:exclusion-radius}

\draftmeta{Chapter~4 (negative-set reconstruction); one sentence in Chapter~7}
{Draft v1 --- decision recorded in log R-036}
{log R-036; \texttt{specs/graphpro\_dataset\_reconstruction.md};
\texttt{specs/sequence\_window\_policy.md}; F-117}

\subsection{Purpose of the radius}

Reconstructed datasets pair every promoter window with a genomic control drawn
from the same assembly. A control is admitted only if it lies outside an
exclusion radius $L_{\mathrm{excl}}$ around every known TSS, on both strands. The radius has
two jobs. First, it keeps known initiation sites out of the negative class.
Second, it prevents a control from sharing sequence with a positive window;
otherwise the same nucleotides would carry opposite labels, the model would
receive contradictory evidence during training, and a grouping procedure would
have to join the two records to avoid leakage across partitions
\citep{teufel2023graphpart,kapoor2023leakage}.

\subsection{Rule}

The window of a domain spans $L_{\mathrm{up}}$ nucleotides upstream of the TSS,
the TSS itself, and $L_{\mathrm{down}}$ nucleotides downstream
(\texttt{specs/sequence\_window\_policy.md}), so that
$L=L_{\mathrm{up}}+1+L_{\mathrm{down}}$. Under a symmetric, strand-agnostic
radius, a control that starts $L_{\mathrm{excl}}+1$ nucleotides from a TSS can
still overlap the upstream part of that TSS's window whenever
$L_{\mathrm{excl}} < L_{\mathrm{up}}$. The radius is therefore defined per
domain as
\begin{equation}
  L_{\mathrm{excl}} = L_{\mathrm{up}} + L_{\mathrm{margin}},
\end{equation}
where $L_{\mathrm{margin}}$ is a margin for imprecision in the TSS position.

\begin{table}[htbp]
  \centering
  \caption{Exclusion radius per taxonomic domain.}
  \begin{tabular}{lrrrr}
    \toprule
    Domain & Window & $L_{\mathrm{up}}$ & $L_{\mathrm{margin}}$ & $L_{\mathrm{excl}}$ \\
    \midrule
    Prokaryota & $-60..{+20}$ (81~nt) & 60 & 40 & 100 \\
    Eukaryota & $-200..{+50}$ (251~nt) & 200 & 50 & 250 \\
    \bottomrule
  \end{tabular}
\end{table}

\subsection{Rationale}

\paragraph{Why the radius must grow with the window.}
The radius of 100~nt was adequate for the original 81-nt window, because it
exceeds $L_{\mathrm{up}}=60$. Once eukaryotic windows were extended to 251~nt, the same
radius would allow a control to share up to 100~nt with the $-200..-101$ region
of a known promoter (log R-036). Tying $L_{\mathrm{excl}}$ to $L_{\mathrm{up}}$ removes this failure for any
window length the protocol may adopt.

\paragraph{Why the eukaryotic margin is wider.}
Bacterial dRNA-seq maps report TSSs at single-nucleotide resolution. Mammalian
promoters, by contrast, frequently initiate transcription over a broad region
rather than at a single position \citep{carninci2006genome,haberle2018eukaryotic},
so a representative TSS such as an EPD record \citep{meylan2020epd} stands for
a distribution of start sites. A larger margin keeps controls away from that
distribution.

\paragraph{Why both strands.}
Promoters can initiate transcription in both directions; in the human genome,
pairs of genes on opposite strands whose start sites lie within 1~kb of each
other are common \citep{trinklein2004abundance}. Excluding only the strand of the
selected TSS would admit controls containing the promoter of the divergent
transcript.

\paragraph{Why the radius is not tuned.}
A smaller radius gives more candidate positions and easier GC matching, and it
produces harder negatives whose genomic context resembles that of promoters.
These are real trade-offs, but the radius defines which sequences count as
non-promoters, and hence the task itself. It is therefore frozen before any
model is trained, never tuned against model outcomes, and any alternative is
run as a separate, versioned sensitivity condition.

\subsection{Feasibility}

The rule does not exhaust the genome even in dense maps. For
\textit{L.~interrogans}, the complete primary TSS table with a 100-nt radius on
both strands leaves 1,029,820 of 4,614,703~bp free, enough for 10,167
non-overlapping 81-nt controls against 2,791 needed; for the 74-kb pAA plasmid
it leaves 41,033~bp, or 453 slots per strand against 124 needed (F-117).

\subsection{Limitations}

The radius protects only against known TSSs. Random genomic controls can still
contain unannotated promoters, so the negative class remains unlabelled rather
than verified, and this uncertainty is reported rather than hidden
\citep{elkan2008pu}.
